%% kapa: exam=1 round=27 year=2016 period=2 form=A subject=accounting subject_ko=회계학 seq=38 number=78
%% source: Q-Net 감정평가사 기출문제 (cst00309) artlSeq=5169402 file=extracted/1008/5169402/제27회 감정평가사 1차 문제/제27회 감정평가사 2교시문제(A형).pdf
%% conversion: kordoc 4.24.51 -> tools/gen1.py (auto); review=proofread-sonnet; flags=none
\begin{question}{78}
(주)감평은 표준원가계산을 적용하고 있으며, 직접노무시간을 기준으로 제조간접원가를 배부하고 있다. 고정제조간접원가 조업도차이는?
\begin{framedtext}
\begin{tabularx}{\linewidth}{@{}Xr@{}}
○ 당기 실제 발생 총제조간접원가 & ￦700,000 \\
○ 실제 직접노무시간 & 70,000시간 \\
○ 기준조업도 & 80,000시간 \\
○ 실제 직접노무시간에서의 제조간접원가 변동예산 & ￦770,000 \\
○ 기준조업도에서의 제조간접원가 예산 & ￦820,000 \\
○ 당기 실제 발생 직접노무원가 & ￦200,000 \\
○ 직접노무원가 임률차이 & ￦25,000(불리) \\
○ 직접노무원가 능률차이 & ￦15,000(유리) \\
○ 표준 직접노무원가 & 시간당 ￦2.5 \\
\end{tabularx}
\end{framedtext}
\begin{choices}
  \item ￦21,000(유리)
  \item ￦21,000(불리)
  \item ￦31,500(유리)
  \item ￦31,500(불리)
  \item ￦52,500(유리)
\end{choices}
\end{question}
\answer{2}
